\makeatletter
\renewcommand {\@degree@string} {Master of Science}
\makeatother

\title{Auto-Adapter: Robot-Specific Driver Synthesis with Independent Closed-Loop Validation}
\author{Yifan Wang}
\department{Department of Computer Science}

\maketitle
\makedeclaration

\begin{abstract}
Large language model (LLM) agents can plan with robot tools only when a
reusable interface and its robot-specific implementation already exist.
Adapting a new robot therefore requires more than successful tool use: the
system must define reusable operations, express testable acceptance
requirements, and synthesise a driver that translates those operations into
robot-specific control without allowing the generated candidate to judge
itself.

This thesis presents Auto-Adapter 2.0, an experiment-grade Direct-MuJoCo
framework for studying these problems. Given public robot facts and
source-backed tasks with reviewed task-level pass standards, Task-Grounded
Capability Design derives five to ten reusable capabilities without a
pre-authored catalogue or task mapping. An implementation-blind Independent
Validation Compiler audits model-authored capability-level criteria against
their source obligations and binds accepted criteria to private executable
cases. A separate trusted Harness observes simulator state and issues final
verdicts. Structured failure reports may support bounded \emph{Repair} within
a run, whereas \emph{Continued Evolution} can expose only human-reviewed
Experience to a later run.

The evaluation comprises three experiments.
\textbf{Experiment 1---Fixed-input LLM backbone comparison.} Seven LLM
backbones are compared across five robot configurations under skeleton-assisted
and from-scratch driver generation. One prior-designed capability interface,
set of pass standards, and private validation suite remains fixed per robot.

\textbf{Experiment 2---Component analysis of Auto-Adapter.} The experiment is
designed to compare reusable-driver synthesis with task-specific code
generation, physics-grounded MuJoCo feedback with \texttt{pytest} feedback, and
bounded Repair with a no-Repair condition. It also evaluates Task-Grounded
Capability Design and the source-preserving formulation, implementation-blind
audit, compilation, and simulator-grounded evaluation of capability-level pass
criteria.

\textbf{Experiment 3---Continued Evolution.} Matched later driver-synthesis
runs compare access to reviewed Experience with the same condition without
Experience.

[RESULT NEEDED: report the principal Experiment~1 synthesis and resource
outcomes, the Experiment~2 component-analysis, capability-design, and
validation-design outcomes, and the matched Experiment~3 Experience effect.]
The findings apply only to the
declared model endpoints, tasks, robot configurations, and Direct-MuJoCo
conditions; they do not establish real-SDK fidelity, hardware safety,
sim-to-real transfer, universal robot or model superiority, autonomous
self-improvement, or generalisation beyond the evaluated conditions.

Within the evaluated scope, the thesis makes three scientific contributions:
\vspace{-.65\baselineskip}
\begin{enumerate}[leftmargin=*,labelsep=.8em,nosep]

  \item AutoAdapter-Bench makes LLM backbone and trusted-skeleton access
  explicit experimental factors while holding the capability interface,
  capability-level pass standards, complete private validation suite, Harness
  rules, and resource budget fixed for each robot. This isolates
  robot-specific driver synthesis from upstream capability design and
  downstream capability-interface use.

  \item Auto-Adapter treats capability abstraction and capability-level
  pass-criteria design as empirical objects rather than hidden human inputs.
  The implementation-blind Independent Validation Compiler audits source
  preservation and binds accepted criteria to private cases, while the trusted
  Harness owns observations and final verdicts. Generation or compilation
  alone is not counted as reliable validation.

  \item A matched evaluation separates within-run Repair from later-run
  Experience reuse. Evolution operates only on a bounded projection of a
  terminal report, and a human disposition is required before a proposal
  becomes Experience. Matched Experience and no-Experience runs test whether
  reviewed evidence changes subsequent driver-synthesis outcomes.

\end{enumerate}

Within the declared Direct-MuJoCo scope, the thesis also makes three bounded,
industry-relevant contributions:
\vspace{-.65\baselineskip}
\begin{enumerate}[leftmargin=*,labelsep=.8em,nosep]

  \item An inspectable path links reviewed task-level standards to
  model-authored capability-level pass criteria, independent compilation, and
  Harness verdicts for the declared tasks and validation cases.

  \item A repeatable fixed-input robot-integration package records the
  simulator assets, morphology and control facts, source-backed tasks,
  capability interface, pass standards, private validation inputs, and
  robot-specific driver used in the experiment.

  \item A review-gated pathway reuses terminal evidence only in later runs.
  Human review, exclusion of private validation content, versioned Experience
  snapshots, and matched no-Experience controls make that reuse inspectable
  without changing the run that produced the evidence.

\end{enumerate}

\end{abstract}

\begin{acknowledgements}
Replace this placeholder with your acknowledgements.
\end{acknowledgements}

\setcounter{tocdepth}{2}
% Setting this higher means you get contents entries for
%  more minor section headers.

\tableofcontents
\listoffigures
\listoftables
